% SPDX-FileCopyrightText: 2026 David Shirvanyants
% SPDX-License-Identifier: CC0-1.0

\chapter{Support Generation}
\label{chap:generate}

Generate converts unsupported layer surfaces into removable triangulated
supports. It snapshots and combines the selected transformed meshes, slices
them, and runs the Detect algorithm. The source mesh, complete layers, and
unsupported polygons then remain read-only while worker threads claim layers,
produce contact points, and construct supports. Reachability of the space
outside the model is calculated concurrently with contact discovery.

\section{Contact Points and Tips}

Each connected unsupported component receives a triangular lattice whose six
nearest neighbors are separated by the configured support spacing. Nodes
outside the component or inside its holes are removed. The algorithm next
finds unsupported perimeter runs which coincide with the source-layer
perimeter. These are exposed model edges; boundaries created only by
subtracting supported area are not exposed. While walking each exposed run,
nearby unused lattice nodes are moved onto the edge at the configured spacing.

An additional pass verifies every unsupported component independently. If no
existing contact in that component can form a valid tip, the center is tested,
followed by interior samples ordered from the center outward. The first valid
sample is added. Thus closely spaced but disconnected islands cannot suppress
one another's contacts, although an island may still remain unsupported when
no geometrically valid tip can be placed.

For each contact, a bounding-volume index finds the nearby source triangle.
Inside/outside probes orient its normal away from the solid. The preferred
axis of a truncated conical tip follows this outward normal from the contact
toward the wide end. If necessary, it is tilted in its vertical plane to
satisfy the critical angle after allowing for the cone's own slope. Samples
along the cone center and circumference reject a tip which enters the model.
Accepted side faces and both ends are triangulated using the configured number
of circumference points.

\section{External Supports}

External supports are preferred because they end on the build platform. The
route search discretizes the model bounds and a surrounding band into a cubic
lattice. At every source layer, solid polygons are expanded by the local
pillar radius, model isolation, and half a cell diagonal as a conservative
grid allowance. These obstacles are stored as compact bit maps.

Cells above collision-free base positions are marked reachable at the lowest
lattice level. Reachability is propagated upward: a cell is accepted when an
unobstructed segment connects it to a reachable cell on the preceding level
and that segment is no shallower than the minimum support angle. Each cell
stores only the direction of its predecessor, which bounds memory use while
retaining a route back to the platform.

Starting from the wide end of a contact tip, the route search connects to this
reachable field and follows predecessor directions down to the base. A direct
probe and a bounded connector search bridge the off-grid tip position. When
the preferred surface-normal tip has no external route, a fallback examines
reachable cells within the tip length and turns the tip toward the nearest
valid one. This sacrifices normal alignment at the contact in order to retain
a collision-free external support. A short transition may ignore only the
extra isolation envelope when the tip necessarily begins within it; physical
pillar clearance is still enforced.

\begin{figure}[htbp]
    \centering
    \begin{tikzpicture}[x=0.75cm,y=0.75cm,>=Latex]
        \fill[gray!35] (0,0) rectangle (11,0.25);
        \draw[thick] (0,0.25) -- (11,0.25)
            node[right] {build platform};
        \fill[gray!25]
            (4.4,7.3) .. controls (6.0,8.1) and (8.2,7.8) ..
            (9.5,6.8) -- (8.8,6.2) .. controls (7.0,7.0) and (5.5,6.8) ..
            (4.4,7.3);
        \draw[thick]
            (4.4,7.3) .. controls (6.0,8.1) and (8.2,7.8) .. (9.5,6.8);
        \draw[thick]
            (4.4,7.3) .. controls (5.5,6.8) and (7.0,7.0) .. (8.8,6.2);
        \draw[very thick, blue!60!black]
            (2.2,0.8) -- (2.2,3.2) -- (3.2,4.2) -- (5.6,5.6);
        \fill[blue!35] (1.5,0.25) rectangle (2.9,0.8);
        \draw[thick] (1.5,0.25) rectangle (2.9,0.8);
        \fill[blue!45]
            (5.25,5.35) -- (5.75,5.85) -- (6.25,6.75) --
            (5.85,6.95) -- cycle;
        \draw[thick]
            (5.25,5.35) -- (5.75,5.85) -- (6.25,6.75) --
            (5.85,6.95) -- cycle;
        \fill (6.05,6.86) circle (0.08);
        \draw[->] (8.1,5.2) -- (6.35,6.55)
            node[pos=0, right] {contact tip};
        \draw[->] (0.8,3.4) -- (2.5,3.4)
            node[pos=0, left] {smoothed pillar};
        \draw[->] (0.5,1.0) -- (1.7,0.65)
            node[pos=0, left] {base};
    \end{tikzpicture}
    \caption{Principal parts of an external support. The centerline may bend
    through reachable space while respecting clearance and angle limits.}
\end{figure}

\section{Internal Fallback and Meshing}

If all external-tip and route choices fail, the generator attempts an internal
support. From the upper tip it advances vertically through lower layers. At
each step, a downward cone limited by the tip height and minimum support angle
is tested against original, unexpanded model contours. The nearest eligible
contour edge becomes the lower contact point. A second tip is placed there
with its narrow end on the model and its wide end toward the pillar.

The lower tip is joined at or above the minimum support angle to a vertical
segment leading to the upper tip. Expanded contours reject a route which
touches the model between its two intended contacts, and insufficient space
between the tips also rejects the support. Internal pillar diameter is
constant and equals the wide diameter of both contact tips.

For either support type, collinear route points are removed and direction
changes are replaced by tangent circular fillets. The base and tip axes are
included in this smoothing. Terminal fillets reduce their radius when a short
base or tip axis cannot accommodate the normal bend; body fillets retain the
local pillar radius. Transported cross-section frames then form a polygonal
tube, while external pillars also receive a tapered radius and cylindrical
platform base.

Every accepted support is a closed triangulated shell. The shells are appended
to one new surface model in the current document without performing a Boolean
union. Contacts for which neither route type succeeds are omitted, counted,
and summarized with the lowest failed Z coordinate. Cooperative cancellation
is checked throughout slicing, contact processing, volume construction, and
route generation; partial results are not installed after cancellation.

\begin{figure}[htbp]
    \centering
    \begin{tikzpicture}[node distance=4mm]
        \node[algstep] (input) {Transformed mesh: slice and detect unsupported polygons};
        \node[algstep, below=of input] (parallel) {In parallel: place per-island contacts and build the external reachability volume};
        \node[algstep, below=of parallel] (tips) {Orient and validate contact tips};
        \node[algdecision, below=7mm of tips] (external) {External route found?};
        \node[algstep, below left=9mm and -4mm of external,
            text width=0.29\textwidth] (platform)
            {Trace to the platform and add a base};
        \node[algstep, below right=9mm and -4mm of external,
            text width=0.29\textwidth] (internal)
            {Try a two-contact internal support};
        \node[algstep, below=19mm of external] (mesh)
            {Smooth centerlines, tessellate closed shells, and append the result model};
        \draw[algarrow] (input) -- (parallel);
        \draw[algarrow] (parallel) -- (tips);
        \draw[algarrow] (tips) -- (external);
        \draw[algarrow] (external) -|
            node[pos=0.25, above] {yes} (platform);
        \draw[algarrow] (external) -|
            node[pos=0.25, above] {no} (internal);
        \draw[algarrow] (platform.south) -- (mesh.north west);
        \draw[algarrow] (internal.south) -- (mesh.north east);
    \end{tikzpicture}
    \caption{Support-generation pipeline.}
\end{figure}
